\section{Mechanism}
\label{sec:mechanism}

The compiled channel is a hypernetwork that amortizes into one forward pass what LoRA
obtains by gradient descent: a trained write function maps a structured belief
representation, through the frozen host's own hidden states, to per-layer conditioning
parameters, in the manner of AdaLN- and FiLM-style conditional normalization
\citep{perez1970film,xu2019adain} and of the hypernetwork-adapter lineage established by
GenerativeAdapter \citep{chen2025generativeadapter} and Text-to-LoRA
\citep{charakorn2025texttolora}.

\subsection{The belief tree}
\label{sec:mechanism-belieftree}

Conditioning content is stored as a \emph{belief tree}: typed hierarchical nodes
\(\{\text{statement}, \text{type} \in \{\text{claim}, \text{argument}, \text{evidence},
\text{experience}, \text{strategy}\}, \text{credence}, \text{edges}
\{\text{supports}, \text{contradicts}\}\}\). Trees are authored by hand for evaluation
topics and produced by the reflection loop for the update-path experiments
(Section~\ref{sec:discrimination}). An early design used numeric credence values on
every node; we found no measurable benefit over a coarse categorical scheme and report
this as a negative result rather than carry the added complexity forward.

\subsection{The modulation}
\label{sec:mechanism-modulation}

The write function conditions query and value projections at a stride-3 subset of
layers (11 of 32), rank 16. For a target layer with base query weight \(W_Q\), the
modulated projection is
\[
  W_Q' = W_Q + \Delta_Q(\phi), \qquad \Delta_Q(\phi) = U_Q(\phi)\,V_Q(\phi)^\top,
\]
with \(U_Q(\phi) \in \mathbb{R}^{d \times 16}\), \(V_Q(\phi) \in \mathbb{R}^{d \times
16}\) produced by the write function from the pooled belief-tree state \(\phi\); the
value projection \(W_V'\) is constructed identically. The compilation is
query-\emph{agnostic}: \(\Delta_Q\) and \(\Delta_V\) do not depend on the current
prompt. Its effect is query-\emph{dependent} anyway, because the low-rank update acts
on the live hidden state produced by whatever the model is currently attending to.

\subsection{The write function}
\label{sec:mechanism-write}

Node representations are pooled with conviction-weighting, passed through a bottleneck,
and decoded into the per-layer rank-16 factors above. Measured directly (60 timed calls,
three topics, \texttt{torch.cuda.synchronize} around the full belief-tree-to-tensors
path), compilation takes 4.83 seconds pooled mean on the evaluation hardware, off the
serving path.

\textbf{Compile, as used here, is lossy.} A trained write function translates the
belief tree into per-layer low-rank tensors, fixed until the store changes, yet
query-dependent in effect because the update acts on the live hidden state. This is a
learned, lossy translation rather than a cached lookup, which is why ``precomputed''
would be the wrong word; the loss is what Section~\ref{sec:boundary} characterizes.
``Compile'' here has no relation to \texttt{torch.compile} or graph compilation.

\subsection{The curated slot}
\label{sec:mechanism-slot}

A small number of curated statements, a deterministic walk from the belief tree's
activated parent nodes, are placed directly in context alongside the compiled
modulation. This slot carries what the bottleneck cannot: specific, low-frequency
content such as named entities, exact figures, or constrained action strings
(Section~\ref{sec:boundary}). Condition \condF{} throughout this paper denotes compiled
modulation plus this curated slot; \condC{} denotes compiled modulation with the slot
empty (zero belief-content context tokens).

\begin{figure}[t]
\centering
\includegraphics[width=\textwidth]{fig_mechanism.png}
\caption{The compile pass and the generation pass. Compilation happens once per belief
change (4.83s measured mean, off the serving path); generation applies the resulting
per-layer tensors on every subsequent forward pass while occupying zero context
positions.}
\label{fig:mechanism}
\end{figure}

\subsection{Serving properties}
\label{sec:mechanism-serving}

A channel choice is also a compute choice, and the two channels this paper compares
scale in opposite directions: context cost grows with belief-token count (mean 90
tokens/belief, range 83--359) and is bounded by the host's context window, while
compiled cost is constant per query, a small same-regime latency addition (0.350s TTFT
vs.\ the bare model's 0.318s), but bounded by a fixed rank-16 budget. The capacity
sweep locates where the compiled budget runs out: per-belief installation holds to
\(k=4\) merged topics and drops at \(k \ge 8\) (Section~\ref{sec:boundary}).
GenerativeAdapter reports a comparable $4\times$ compute-and-storage reduction against
full-conversation prompting on its own personalization benchmark
\citep{chen2025generativeadapter}; our addition is that the channel which saves the
compute is also the channel that holds the stance, which that comparison was not built
to measure.\footnote{Full cost analysis and write-function architecture in
Appendix~\ref{app:serving}--\ref{app:writefn}.}
